% Options for packages loaded elsewhere
\PassOptionsToPackage{unicode}{hyperref}
\PassOptionsToPackage{hyphens}{url}
\documentclass[
  ignorenonframetext,
  aspectratio=169,
]{beamer}
\newif\ifbibliography
\usepackage{pgfpages}
\setbeamertemplate{caption}[numbered]
\setbeamertemplate{caption label separator}{: }
\setbeamercolor{caption name}{fg=normal text.fg}
\beamertemplatenavigationsymbolsempty
% remove section numbering
\setbeamertemplate{part page}{
  \centering
  \begin{beamercolorbox}[sep=16pt,center]{part title}
    \usebeamerfont{part title}\insertpart\par
  \end{beamercolorbox}
}
\setbeamertemplate{section page}{
  \centering
  \begin{beamercolorbox}[sep=12pt,center]{section title}
    \usebeamerfont{section title}\insertsection\par
  \end{beamercolorbox}
}
\setbeamertemplate{subsection page}{
  \centering
  \begin{beamercolorbox}[sep=8pt,center]{subsection title}
    \usebeamerfont{subsection title}\insertsubsection\par
  \end{beamercolorbox}
}
% Prevent slide breaks in the middle of a paragraph
\widowpenalties 1 10000
\raggedbottom
\AtBeginPart{
  \frame{\partpage}
}
\AtBeginSection{
  \ifbibliography
  \else
    \frame{\sectionpage}
  \fi
}
\AtBeginSubsection{
  \frame{\subsectionpage}
}
\usepackage{iftex}
\ifPDFTeX
  \usepackage[T1]{fontenc}
  \usepackage[utf8]{inputenc}
  \usepackage{textcomp} % provide euro and other symbols
\else % if luatex or xetex
  \usepackage{unicode-math} % this also loads fontspec
  \defaultfontfeatures{Scale=MatchLowercase}
  \defaultfontfeatures[\rmfamily]{Ligatures=TeX,Scale=1}
\fi
\usepackage{lmodern}
\ifPDFTeX\else
  % xetex/luatex font selection
\fi
% Use upquote if available, for straight quotes in verbatim environments
\IfFileExists{upquote.sty}{\usepackage{upquote}}{}
\IfFileExists{microtype.sty}{% use microtype if available
  \usepackage[]{microtype}
  \UseMicrotypeSet[protrusion]{basicmath} % disable protrusion for tt fonts
}{}
\makeatletter
\@ifundefined{KOMAClassName}{% if non-KOMA class
  \IfFileExists{parskip.sty}{%
    \usepackage{parskip}
  }{% else
    \setlength{\parindent}{0pt}
    \setlength{\parskip}{6pt plus 2pt minus 1pt}}
}{% if KOMA class
  \KOMAoptions{parskip=half}}
\makeatother
\usepackage{longtable,booktabs,array}
\usepackage{caption}
\captionsetup[table]{skip=6pt}
\newcounter{none} % for unnumbered tables
\usepackage{calc} % for calculating minipage widths
\usepackage{caption}
% Make caption package work with longtable
\makeatletter
\def\fnum@table{\tablename~\thetable}
\makeatother
\setlength{\emergencystretch}{3em} % prevent overfull lines
\providecommand{\tightlist}{%
  \setlength{\itemsep}{0pt}\setlength{\parskip}{0pt}}
% Auto-generated by infrastructure/rendering/_pdf_latex_helpers.py::extract_math_font_preamble.
% Minimal header passed to Pandoc via -H so Beamer slides pick up the same math font
% as the combined-PDF path. Adjust the manuscript's preamble.md to change the font globally.
\usepackage{unicode-math}
\setmathfont{latinmodern-math.otf}

% Manuscript macro/package fallbacks (newcommand -> providecommand).
% Core mathematics
\usepackage{amsmath}
\usepackage{amssymb}
% Document layout
% Tables
\usepackage{booktabs}
% Caption styling (booktabs already loaded above). Small caption font with a
% bold label ("Figure 1", "Table 2") gives the math/figure-heavy layout a
% consistent, journal-like caption rhythm without fighting pandoc-crossref —
% pandoc emits standard \caption{}, and caption/captionsetup only restyle it.
% ── Theorem environments ─────────────────────────────────────────────────
% amsthm is loaded above. The FedGVI<->Friston bridge is stated as a small
% number of theorems/definitions/lemmas/propositions/corollaries; sharing one
% counter (definition/lemma/proposition/corollary all step the `theorem`
% counter) gives a single monotone numbering 1, 2, 3, ... across all kinds, so
% authors can reference them as Theorem 1, Definition 2, Corollary 3 without a
% per-kind counter drifting out of order. Plain style for theorem-like results,
% definition style (upright body) for definitions.
% (algorithm2e intentionally not loaded: this manuscript states results as
% numbered amsthm Theorems/Definitions, not algorithm2e pseudocode.)
% Code listings
\usepackage{listings}
% Typography and formatting
% (siunitx intentionally not loaded: no \SI/\num units appear in this manuscript.)
% Cross-references and citations
% (cleveref and titlesec intentionally not loaded: unavailable in this TeX tree
% and not required. Cross-references resolve through pandoc-crossref's own
% \ref-style names — no \cref appears — and section heading styling falls back
% to the pandoc/LaTeX defaults.)
% ── TeX Live Basic-compatible mono font for code listings ────────────
% Latin Modern Mono is bundled with TeX Live Basic and keeps the manuscript
% renderable on minimal TeX installations. Mathematical Unicode coverage is
% provided separately by Latin Modern Math below, so the code font does not
% need a private JuliaMono installation.
%
% Use the TeX font filename rather than the system-family name: the minimal
% TeX Live fontconfig database may not register the OpenType family even though
% kpsewhich can resolve the bundled files.
% Math font for unicode-math: Latin Modern Math (TeX Live) has full BMP
% coverage including U+2223 (\mid), U+226A/226B (\ll/\gg), and the Greek/
% blackboard letters used in equations. Without an explicit \setmathfont,
% unicode-math falls back to lmroman text font which lacks several glyphs
% and emits "Missing character" warnings on every \mid in math mode.

% Slide layout defaults for warning-clean scientific decks.
\usepackage{etoolbox}
\IfFileExists{xurl.sty}{\usepackage{xurl}}{}
\IfFileExists{seqsplit.sty}{\usepackage{seqsplit}}{\newcommand{\seqsplit}[1]{#1}}
\protected\def\breakseq#1{\seqsplit{#1}}
\protected\def\breaktt#1{\begingroup\ttfamily\seqsplit{#1}\endgroup}
\setlength{\emergencystretch}{6em}
\tolerance=5000
\hbadness=10000
\hfuzz=1pt
\setlength{\tabcolsep}{2pt}
\AtBeginEnvironment{longtable}{\tiny\renewcommand{\arraystretch}{0.86}\setlength{\tabcolsep}{1pt}}
\AtBeginEnvironment{tabular}{\tiny\renewcommand{\arraystretch}{0.86}\setlength{\tabcolsep}{1pt}}
\AtBeginEnvironment{equation}{\tiny}
\AtBeginEnvironment{equation*}{\tiny}
\AtBeginEnvironment{align}{\tiny}
\AtBeginEnvironment{align*}{\tiny}
\AtBeginEnvironment{itemize}{\footnotesize}
\AtBeginEnvironment{enumerate}{\footnotesize}
\AtBeginEnvironment{description}{\footnotesize}
\setbeamerfont{caption}{size=\tiny}
\setbeamerfont{caption name}{size=\tiny}
\setbeamerfont{normal text}{size=\small}
\setbeamerfont{frametitle}{size=\small}
\setbeamerfont{section title}{size=\footnotesize}
\setbeamerfont{subsection title}{size=\footnotesize}
\setbeamertemplate{section page}{%
  \centering
  \begin{beamercolorbox}[sep=12pt,center,wd=\paperwidth]{section title}
    \parbox{0.86\paperwidth}{\centering\usebeamerfont{section title}\insertsection\par}
  \end{beamercolorbox}
}
\setbeamertemplate{subsection page}{%
  \centering
  \begin{beamercolorbox}[sep=8pt,center,wd=\paperwidth]{subsection title}
    \parbox{0.86\paperwidth}{\centering\usebeamerfont{subsection title}\insertsubsection\par}
  \end{beamercolorbox}
}
\setlength{\abovecaptionskip}{2pt}
\setlength{\belowcaptionskip}{0pt}

% Natbib and cross-reference fallbacks — slides don't load natbib
% or cleveref, but combined-PDF manuscript prose may emit these
% commands. The fallback renders citations as a bracketed key list
% and cross-references as detokenized labels so slides don't fail on
% undefined control sequences or raw underscores. \providecommand is
% a no-op if packages load later.
\providecommand{\citep}[1]{[#1]}
\providecommand{\citet}[1]{#1}
\providecommand{\citealp}[1]{#1}
\providecommand{\citeauthor}[1]{#1}
\providecommand{\citeyear}[1]{#1}
\providecommand{\cref}[1]{\texttt{\detokenize{#1}}}
\providecommand{\Cref}[1]{\texttt{\detokenize{#1}}}
\providecommand{\crefrange}[2]{\texttt{\detokenize{#1}}--\texttt{\detokenize{#2}}}
\providecommand{\Crefrange}[2]{\texttt{\detokenize{#1}}--\texttt{\detokenize{#2}}}
\providecommand{\paragraph}[1]{\textbf{#1}\ }

% Opt-in accessible presentation profile.
\setbeamerfont{normal text}{size*={20pt}{24pt}}
\setbeamerfont{frametitle}{size*={28pt}{32pt}}
\setbeamerfont{section title}{size*={28pt}{32pt}}
\setbeamerfont{subsection title}{size*={28pt}{32pt}}
\setbeamerfont{caption}{size*={16pt}{19pt}}
\setbeamerfont{caption name}{size*={16pt}{19pt}}
\setbeamerfont{itemize/enumerate body}{size*={20pt}{24pt}}
\setbeamerfont{itemize/enumerate subbody}{size*={20pt}{24pt}}
\setbeamerfont{itemize/enumerate subsubbody}{size*={20pt}{24pt}}
\setbeamerfont{itemize item}{size*={20pt}{24pt}}
\setbeamerfont{itemize subitem}{size*={20pt}{24pt}}
\setbeamerfont{itemize subsubitem}{size*={20pt}{24pt}}
\setbeamerfont{description body}{size*={20pt}{24pt}}
\setbeamerfont{description item}{size*={20pt}{24pt}}
\setbeamerfont{quote}{size*={20pt}{24pt}}
\setbeamerfont{quotation}{size*={20pt}{24pt}}
\AtBeginEnvironment{longtable}{\fontsize{20pt}{24pt}\selectfont}
\AtBeginEnvironment{tabular}{\fontsize{20pt}{24pt}\selectfont}
\AtBeginEnvironment{equation}{\fontsize{20pt}{24pt}\selectfont}
\AtBeginEnvironment{equation*}{\fontsize{20pt}{24pt}\selectfont}
\AtBeginEnvironment{align}{\fontsize{20pt}{24pt}\selectfont}
\AtBeginEnvironment{align*}{\fontsize{20pt}{24pt}\selectfont}
\AtBeginEnvironment{itemize}{\fontsize{20pt}{24pt}\selectfont}
\AtBeginEnvironment{enumerate}{\fontsize{20pt}{24pt}\selectfont}
\AtBeginEnvironment{description}{\fontsize{20pt}{24pt}\selectfont}
\AtBeginDocument{\usebeamerfont{normal text}}
\newlength{\AFSlideTableWidth}
% The fixed footer reservation is calibrated with the semantic seven-line regular-body budget.
\setbeamertemplate{footline}{%
  \leavevmode\vbox{%
    \hbox{%
      \begin{beamercolorbox}[wd=\paperwidth,ht=16pt,dp=3pt,center]{author in head/foot}%
        \fontsize{16pt}{19pt}\selectfont Untagged PDF derivative \textbar\ \href{../web/index.html}{HTML reader}%
      \end{beamercolorbox}%
    }%
    \vskip6pt%
  }%
}
% Accessible footer reserved height: 25pt.

\newtheorem{proposition}{Proposition}
\newtheorem{hypothesis}{Hypothesis}
\newtheorem{remark}[theorem]{Remark}
\newtheorem{axiom}{Axiom}
\newtheorem{property}{Property}
\usepackage{bookmark}
\IfFileExists{xurl.sty}{\usepackage{xurl}}{} % add URL line breaks if available
\urlstyle{same}
% fallback for those not using the hyperref driver hyperxmp:
\makeatletter
\@ifundefined{xmpquote}{\newcommand{\xmpquote}[1]{#1}}{}
\makeatother
\hypersetup{
  hidelinks,
  pdfcreator={LaTeX via pandoc}}

\author{\texorpdfstring{}{}}
\date{}

\begin{document}

\begin{frame}{Generalized Bayes: the route back to standard Bayes}
\protect\phantomsection\label{sec:method-genbayes}
The inference engine FedGVI federates is the generalized (Gibbs)
posterior (Bissiri et al. 2016; Jiang and Tanner 2008; Jewson et al.
2018; Knoblauch et al. 2022),
\end{frame}

\begin{frame}{Generalized Bayes (part 2)}
which trades the likelihood for a loss \(L\) and the KL regularizer for
a general divergence \(\mathcal D\):
\end{frame}

\begin{frame}{Generalized Bayes (part 3)}
\begin{equation}\tag{2}\protect\phantomsection\label{eq:gen-bayes}{
\begin{aligned}
q_n^\ast(s)
&= \arg\min_{q_n}
\left\{
\begin{aligned}
&\mathbb{E}_{q_n}\!\Big[\textstyle\sum_i L(s; o_i)\Big]\\
&\quad + \tfrac{1}{\tau}\,\mathcal D\!\big(q_n \,\|\, \pi_0\big)
\end{aligned}
\right\}.
\end{aligned}
}\end{equation}
\end{frame}

\begin{frame}{Generalized Bayes (part 4)}
with prior \(\pi_0\), learning rate \(\tau\), and regularizing
divergence \(\mathcal D\). The learning rate is part of the inferential
specification, not a cosmetic constant;
\end{frame}

\begin{frame}{Generalized Bayes (part 5)}
coarsened-posterior and safe-Bayes work show why calibration of that
temperature matters under misspecification (Miller and Dunson 2018;
Grünwald 2012), where ordinary Bayes concentrates around a KL
pseudo-truth rather than literal truth when the model family is wrong
(Kleijn and Vaart 2012).
\end{frame}

\begin{frame}{Generalized Bayes (part 6)}
We name the object eq.~2 defines.
\end{frame}

\begin{frame}{Generalized Bayes (part 7)}
\begin{definition}[Generalized-(Gibbs)-Bayes posterior]\label{def:generalized-bayes}
Given loss $L$, prior $\pi_0$, learning rate $\tau>0$, and divergence
$\mathcal D$, the generalized-Bayes posterior $q_n^\ast$ minimizes
(2). When $\mathcal D=\mathrm{KL}$, its exact categorical
minimizer is the tempered softmax in (3).
\end{definition}
\end{frame}

\begin{frame}{Generalized Bayes (part 8)}
\begin{equation}\tag{3}\protect\phantomsection\label{eq:tempered-softmax}{
q_n^\ast(s) \;\propto\; \pi_0(s)\,\exp\!\big(-\tau \textstyle\sum_i L(s; o_i)\big),
}\end{equation}
\end{frame}

\begin{frame}[fragile]{Generalized Bayes (part 9)}
The implementation surface is the \texttt{generalized\_posterior}
function in the \texttt{generalized\_bayes} module.
\end{frame}

\begin{frame}{Generalized Bayes (part 10)}
The tempered softmax of eq.~3, stated in the
definition above, is not an approximation: it is the exact closed-form
minimizer of eq.~2 when the regularizer is the KL
divergence, because the categorical support is finite and the objective
is strictly convex in \(q\).
\end{frame}

\begin{frame}{Generalized Bayes (part 11)}
The recovery to standard Bayes follows by choosing the loss.
\end{frame}

\begin{frame}{Generalized Bayes (part 12)}
With \(L=\mathrm{NLL}\), \(\mathrm{NLL}(p, o) = -\log p(o)\), the
exponential in that tempered softmax becomes a product of likelihoods
and the minimizer is \emph{exactly} standard Bayes;
eq.~11 in sec.~5.8 states
that corner,
\end{frame}

\begin{frame}{Generalized Bayes (part 13)}
and Corollary 6 there pins it to the
closed-form prior-times-likelihood product.
\end{frame}

\begin{frame}[fragile]{Generalized Bayes (part 14)}
The largest observed discrepancy between \texttt{generalized\_posterior}
in this regime and the analytic Bayes posterior is 5.55e-17, reported in
sec.~7.1 --- exact to machine precision (a
maximum deviation of about one ULP), not merely close.
\end{frame}

\begin{frame}{Generalized Bayes (part 15)}
FedGVI computes each client update against a \emph{cavity} rather than
the full posterior, so a contributing agent does not double-count its
own previous message. We name that operation.
\end{frame}

\begin{frame}{Generalized Bayes (part 16)}
\begin{definition}[Cavity / PVI factor update]\label{def:cavity}
The cavity is the normalized removal of agent $n$'s site factor from the colony
posterior in natural-parameter space, as in (4). PVI replaces the
site factor and recombines it with that cavity according to
(5).
\end{definition}
\end{frame}

\begin{frame}{Generalized Bayes (part 17)}
\begin{equation}\tag{4}\protect\phantomsection\label{eq:cavity}{
\begin{aligned}
q_{-n}(s)
&=
\frac{q(s)/t_n(s)}{\sum_{s'} q(s')/t_n(s')}
\\
&=
\operatorname{softmax}\!\big(\log q(s)-\log t_n(s)\big),
\end{aligned}
}\end{equation}
\end{frame}

\begin{frame}{Generalized Bayes (part 18)}
The final expression makes normalization explicit. Taking a cavity and
re-multiplying the original site factor restores the global posterior:
\end{frame}

\begin{frame}{Generalized Bayes (part 19)}
\begin{equation}\tag{5}\protect\phantomsection\label{eq:factor-replacement}{
q(s)=\frac{q_{-n}(s)t_n(s)}{\sum_{s'}q_{-n}(s')t_n(s')},
}\end{equation}
\end{frame}

\begin{frame}[fragile]{Generalized Bayes (part 20)}
This is the recombination identity satisfied by
\texttt{generalized\_bayes.cavity} and
\texttt{generalized\_bayes.update\_factor}.

The numbered recombination identity is eq.~5.
\end{frame}

\begin{frame}{Generalized Bayes (part 21)}
The cavity of eq.~4 is the discrete analogue of the
expectation- propagation / partitioned-VI cavity used outside active
inference (Ashman et al. 2022; Bui et al. 2018),
\end{frame}

\begin{frame}{Generalized Bayes (part 22)}
imported here so that the per-agent generalized-Bayes update of
eq.~2 is computed against the colony belief with the
agent's own contribution removed --- exactly the sensory-attenuation
discipline the belief-sharing round of
sec.~5.9 requires.
\end{frame}

\begin{frame}{Generalized Bayes (part 23)}
What remains unspecified in eq.~2 are its two
ingredients --- the divergence \(\mathcal D\) and the loss \(L\) ---
whose robust members and standard-Bayes limits
sec.~5.6 develops next;
\end{frame}

\begin{frame}{Generalized Bayes (part 24)}
the aggregation identity (sec.~5.8) then
federates the resulting per-agent posteriors.
\end{frame}

\begin{frame}{Generalized Bayes (part 25)}
The authoritative notation supplement makes the same normalization and
recombination contract explicit in eq.~36 and
eq.~37; those equations govern the
symbols used by the implementation and all later supplements.
\end{frame}

\begin{frame}{Conjugate likelihood learning for the shared model}
\protect\phantomsection\label{sec:method-learning}
Active-inference agents learn the parameters of their generative model,
not just plan with them (Smith et al. 2022; Friston et al. 2024).
\end{frame}

\begin{frame}{Conjugate likelihood\ldots{} (part 2)}
The likelihood matrix \(A\) carries a Dirichlet prior with concentration
\(a\) over each column, updated conjugately by accumulating
observation-state co-occurrence counts (eq.~18),
\end{frame}

\begin{frame}{Conjugate likelihood\ldots{} (part 3)}
giving the column-normalized expected likelihood. The update of
eq.~18 is driven by the expected sufficient
statistics under the data-generating model,

so as the concentrations accumulate \(\mathbb{E}[A]\) converges to the
true likelihood.
\end{frame}

\begin{frame}{Conjugate likelihood\ldots{} (part 4)}
Convergence is measured by the per-column KL divergence summed over
hidden states, which decreases monotonically toward the standard-Bayes
fixed point; sec.~7.3 reports the learning curve,
where the KL falls from 3.4231 to 0.0027 across 24 count batches.
\end{frame}

\begin{frame}{Conjugate likelihood\ldots{} (part 5)}
A forgetting hyperprior optionally decays the running mass toward an
asymptote so the agent does not become infinitely confident; with the
hyperprior disabled the classical unbounded accumulation of
eq.~18 is recovered.
\end{frame}

\begin{frame}[fragile]{Conjugate likelihood\ldots{} (part 6)}
The implementation is \texttt{learn\_likelihood} in the
\texttt{dirichlet\_learning} module.
\end{frame}

\begin{frame}[fragile]{Bayesian model reduction for structure
comparison}
\protect\phantomsection\label{sec:method-bmr}
Structure learning in the active-inference frame proceeds by Bayesian
model reduction (BMR): given a full model with Dirichlet posterior
\texttt{post} under prior \texttt{prior},
\end{frame}

\begin{frame}{Bayesian model reduction\ldots{} (part 2)}
the change in negative variational free energy from swapping in a
\emph{reduced} prior --- for example one that prunes a redundant column
toward zero ---
\end{frame}

\begin{frame}{Bayesian model reduction\ldots{} (part 3)}
is available in closed form without re-running inference (Friston and
Penny 2011; Smith et al. 2022).
\end{frame}

\begin{frame}[fragile]{Bayesian model reduction\ldots{} (part 4)}
Because the likelihood is shared, the reduced posterior is
\texttt{post\ +\ reduced\_prior\ -\ prior}, and the free-energy
difference is a difference of log multivariate Beta functions
(eq.~20), where
\(\ln B(a) = \sum_k \ln\Gamma(a_k) - \ln\Gamma(\sum_k a_k)\) is the log
Dirichlet normalizer.
\end{frame}

\begin{frame}{Bayesian model reduction\ldots{} (part 5)}
A positive \(\Delta F\) in eq.~20 means the reduced
model has more evidence --- the pruned structure was redundant and
should be adopted; a negative \(\Delta F\) means the reduction destroyed
something the data support. When the reduced prior equals the prior the
score is identically zero,
\end{frame}

\begin{frame}[fragile]{Bayesian model reduction\ldots{} (part 6)}
the no-reduction fixed point.

sec.~7.4 reports \(\Delta F = 3.68\) for a
redundant reduction (accepted) against \(\Delta F = -27.67\) for a
supported one (rejected). The implementation is
\texttt{bayesian\_model\_reduction.reduce}.
\end{frame}

\end{document}
